% ============ §23 Flatness and fibres (原书 pp.178-185) ============
\section{Flatness and fibres}\label{sec:23}\index{flatness and fibres, section 23}\index{fibre ring, section 23}\index{normal ring, criterion of Serre, section 23}\index{regularity and flatness, section 23}

Let $(A, m)$ and $(B, n)$ be Noetherian local rings, and $\varphi\colon A \to B$ a local
homomorphism. We set $F = B \otimes_A k(m) = B/mB$ for the \emph{fibre ring}\index{fibre ring} of
$\varphi$ over $m$. If $B$ is flat over $A$ then according to Theorem~\ref{thm:15.1}, we
have
\[
  (*) \quad \dim B = \dim A + \dim F.
\]

As the following shows, under certain conditions the converse holds.

\begin{theorem}\label{thm:23.1}
Let $A$, $B$ and $F$ be as above. If $A$ is a regular local ring, $B$ is
Cohen-Macaulay, and $\dim B = \dim A + \dim F$ then $B$ is flat over $A$.
\end{theorem}

\begin{proof}
We use induction on $\dim A$. If $\dim A = 0$ then $A$ is a field, and we are
done. If $\dim A > 0$, take $x \in m - m^2$ and set $A' = A/xA$ and $B' = B/xB$. Then
\[
  \dim B' \le \dim A' + \dim F = \dim A - 1 + \dim F = \dim B - 1,
\]
and using a system of parameters of $B'$ one sees that $\dim B' \ge \dim B - 1$, so
that
\[
  \dim B' = \dim A' + \dim F = \dim B - 1.
\]
One sees easily from this that $x$ is $B$-regular and $B'$ is a CM ring. Hence by
induction $B'$ is flat over $A'$. Thus $\Tor^{A'}_1(A/m, B') = 0$; moreover, $x$ is
both $A$-regular and $B$-regular, so that $\Tor^A_1(A/m, B') = \Tor^A_1(A/m, B)$.
Therefore by Theorem~\hyperref[thm:22.3]{22.3}, $B$ is flat over $A$.
\end{proof}

We give a translation of the above theorem into algebraic geometry for the sake
of application. (The language is that of modern algebraic geometry, see for
example [Ha], Ch.~2.)

\begin{corollary*}
Let $k$ be a field, $X$ and $Y$ irreducible algebraic $k$-schemes, and
$f\colon Y \to X$ a morphism. Set $\dim X = n$, $\dim Y = m$, and suppose that the
following conditions hold: (1) $X$ is regular; (2) $Y$ is Cohen-Macaulay; (3) $f$
takes closed points of $Y$ into closed points of $X$ (this holds for example if
$f$ is proper); (4) for every closed point $x \in X$ the fibre $f^{-1}(x)$ is
$(m - n)$-dimensional (or empty). Then $f$ is flat.
\end{corollary*}

\begin{proof}
Let $y \in Y$ be a closed point, and set $x = f(y)$, $A = \mathcal{O}_{X,x}$, and
$B = \mathcal{O}_{Y,y}$. We have $\dim A = n$, $\dim B = m$, and since by
Theorem~\ref{thm:15.1} $\dim B/m_AB \le \dim B - \dim A = m - n$, we get
$\dim B/m_AB = m - n$ from (4). Therefore by the above theorem $B$ is flat over $A$,
and this is what was required to prove.
\end{proof}

\begin{theorem}\label{thm:23.2}
Let $\varphi\colon A \to B$ be a homomorphism of Noetherian rings, and let $E$ be an
$A$-module and $G$ a $B$-module. Suppose that $G$ is flat over $A$; then we have
the following:
\begin{enumerate}[(i)]
  \item if $p \in \Spec A$ and $G/pG \neq 0$ then $\Ass_A(G/pG) = \{p\}$; moreover, if
        $P \in \Ass_B(G/pG)$ then $P \cap A = p$;
  \item $\Ass_B(E \otimes_A G) = \bigcup_{p \in \Ass_A(E)} \Ass_B(G/pG)$.
\end{enumerate}
\end{theorem}

\begin{proof}
(i) $G/pG = G \otimes_A (A/p)$ is flat over $A/p$, and $A/p$ is an integral domain,
so that any non-zero element of $A/p$ is $G/pG$-regular (see Ex.~\ref{ex:7.5-ex}). In
other words, the elements of $A - p$ are $G/pG$-regular. This gives
$\Ass_A(G/pG) = \{p\}$. Also, if $P \in \Ass_B(G/pG)$ then there exists $t \in G/pG$
such that $\ann_B(t) = P$, and then $P \cap A = \ann_A(t) \in \Ass_A(G/pG) = \{p\}$.

(ii) If $p \in \Ass_A(E)$ then there is an exact sequence of the form
$0 \to A/p \to E$, and since $G$ is flat the sequence $0 \to G/pG \to E \otimes_A G$
is also exact; thus
\[
  \Ass_B(G/pG) \subset \Ass_B(E \otimes_A G).
\]
Conversely, if $P \in \Ass_B(E \otimes_A G)$ then there is an $e \in E \otimes_A G$ such
that $\ann_B(e) = P$. We write $e = \sum x_i \otimes y_i$ with $x_i \in E$ and $y_i \in G$,
and set $E' = \sum Ax_i$; then by flatness of $G$, we can view $E' \otimes_A G$ as a
submodule $E' \otimes_A G \subset E \otimes_A G$. Since $e \in E' \otimes_A G$ we have
$P \in \Ass_B(E' \otimes_A G)$. Now $E'$ is a finite $A$-module, so that we can
choose a shortest primary decomposition of $0$ in $E'$, say
$0 = Q_1 \cap \cdots \cap Q_r$. Since $E'$ can be embedded in $\bigoplus (E'/Q_i)$, if we set
$E_i = E'/Q_i$ then
\[
  \Ass_B(E' \otimes_A G) \subset \bigcup_i \Ass_B(E_i \otimes_A G),
\]
and therefore $P \in \Ass_B(E_i \otimes_A G)$ for some $i$. This $E_i$ is a finite
$A$-module having just one associated prime, say $p$. We have
$p \in \Ass_A(E') \subset \Ass_A(E)$. For large enough $\nu$ we get $p^\nu E_i = 0$, so
that $p^\nu(E_i \otimes_A G) = 0$, and thus $p \subset P \cap A$. Moreover, an
element of $A - p$ is $E_i$-regular, and hence also $E_i \otimes_A G$-regular, so that
finally $p = P \cap A$. Now choose a chain of submodules of $E_i$,
\[
  E_i = E_0 \supset E_1 \supset \cdots \supset E_r = 0
\]
such that $E_j/E_{j+1} \cong A/p_j$ with $p_j \in \Spec A$. Then also
\[
  E_i \otimes_A G \supset E_1 \otimes_A G \supset \cdots \supset E_r \otimes_A G = 0,
\]
with
\[
  (E_j \otimes_A G)/(E_{j+1} \otimes_A G) \cong (A/p_j) \otimes_A G = G/p_jG,
\]
so that $\Ass_B(E_i \otimes_A G) \subset \bigcup_j \Ass_B(G/p_jG)$. Therefore
$P \in \Ass_B(G/p_jG)$ for some $j$, but by (i), $P \cap A = p_j$, so that $p_j = p$
and $P \in \Ass_B(G/pG)$.
\end{proof}

\begin{theorem}\label{thm:23.3}
Let $(A, m, k)$ and $(B, n, k')$ be Noetherian local rings, and
$\varphi\colon A \to B$ a local homomorphism. Let $M$ be a finite $A$-module, $N$ a
finite $B$-module, and assume that $N$ is flat over $A$. Then
\[
  \depth_B(M \otimes_A N) = \depth_A M + \depth_B(N/mN).
\]
\end{theorem}

\begin{proof}
Let $x_1, \ldots, x_r \in m$ be a maximal $M$-sequence, and $y_1, \ldots, y_s \in n$ a
maximal $N/mN$-sequence. Writing $x_i'$ for the images of $x_i$ in $B$, let us
prove that $x_1', \ldots, x_r', y_1, \ldots, y_s$ is a maximal $M \otimes_A N$-sequence. Now
$x_1', \ldots, x_r'$ is an $M \otimes_A N$-sequence, and if we set
$M_1 = M/\sum x_iM$ then
\[
  m \in \Ass_A(M_1), \quad \text{and} \quad (M \otimes_A N) \Big/ \sum_{i=1}^r x_i(M \otimes_A N) = M_1 \otimes_A N.
\]
Moreover, by the corollary of Theorem~\hyperref[thm:22.5]{22.5}, $y_1$ is $N$-regular, and
$N_1 = N/y_1N$ is flat over $A$, so that from the exact sequence
$0 \to N \xrightarrow{y_1} N \to N_1 \to 0$ we get the exact sequence
$0 \to M_1 \otimes_A N \xrightarrow{y_1} M_1 \otimes_A N \to M_1 \otimes_A N_1 \to 0$.
Proceeding in the same way we see that $y_1, \ldots, y_s$ is an
$M_1 \otimes_A N$-sequence. After this we need only prove that the $B$-module
\[
  (M \otimes_A N) \Big/ \Big( \sum x_i(M \otimes_A N) + \sum y_j(M \otimes_A N) \Big) = M_1 \otimes_A N_s
\]
has depth $0$, that is $n \in \Ass_B(M_1 \otimes_A N_s)$; however,
$m \in \Ass_A(M_1)$ and $n \in \Ass_B(N_s/mN_s)$, so that this follows at once from
the previous theorem.
\end{proof}

\begin{corollary*}
Let $A \to B$ be a local homomorphism of Noetherian rings as in the theorem, and
set $F = B/mB$. Assume that $B$ is flat over $A$. Then
\begin{enumerate}[(i)]
  \item $\depth B = \depth A + \depth F$;
  \item $B$ is CM $\iff$ $A$ and $F$ are both CM.
\end{enumerate}
\end{corollary*}

\begin{proof}
(i) is the case $M = A$, $N = B$ of the theorem. From (i) and (*) we have
\[
  \dim B - \depth B = (\dim A - \depth A) + (\dim F - \depth F)
\]
and in view of $\dim A \ge \depth A$ and $\dim F \ge \depth F$, (ii) is clear.
\end{proof}

\begin{theorem}\label{thm:23.4}
Let $A \to B$ be a local homomorphism of Noetherian local rings, set $m = \rad(A)$
and $F = B/mB$. We assume that $B$ is flat over $A$; then
\[
  B \text{ is Gorenstein} \iff A \text{ and } F \text{ are both Gorenstein}.
\]
\end{theorem}

\begin{proof}
(K.~Watanabe [1]). By the corollary just proved, we can assume that $A$, $B$ and
$F$ are CM. Set $\dim A = r$ and $\dim F = s$, and let $\{x_1, \ldots, x_r\}$ be a
system of parameters of $A$, and $\{y_1, \ldots, y_s\}$ a subset of $B$ which reduces
to a system of parameters of $F$ modulo $mB$. Then as we have seen in the proof
of Theorem~\hyperref[thm:23.3]{3}, $(x_1, \ldots, x_r, y_1, \ldots, y_s)$ is a $B$-sequence, and
therefore a system of parameters of $B$, and $\bar{B} = B/(x, y)B$ is flat over
$\bar{A} = A/(x)A$. Thus replacing $A$ and $B$ by $\bar{A}$ and $\bar{B}$, we can reduce to
the case $\dim A = \dim B = 0$. Now in general, a zero-dimensional local ring
$(R, M)$ is Gorenstein if and only if $\Hom_R(R/M, R) = (0 : M)_R$ is isomorphic to
$R/M$. Now set
\[
  \rad(B) = n, \quad \rad(F) = n/mB = \bar{n} \quad \text{and} \quad (0 : m)_A = I.
\]
Then $I$ is of the form $I \cong (A/m)^t$ for some $t$, and
$(0 : mB)_B = IB \cong (A/m)^t \otimes_A B = F^t$. Furthermore, we have
$(0 : n)_B = (0 : \bar{n})_{IB} \cong ((0 : \bar{n})_F)^t$, and hence if we set
$(0 : \bar{n})_F \cong (F/\bar{n})^u = (B/n)^u$ then $(0 : n)_B \cong (B/n)^{tu}$.
Therefore
\[
  B \text{ is Gorenstein} \iff tu = 1 \iff t = u = 1 \iff A \text{ and } F \text{ are Gorenstein}.
\]
\end{proof}

\begin{theorem}\label{thm:23.5}
If $A$ is Gorenstein then so are $A[X]$ and $A[[X]]$.
\end{theorem}

\begin{proof}
We write $B$ for either of $A[X]$ or $A[[X]]$, so that $B$ is flat over $A$. For any
maximal ideal $M$ of $B$ we set $M \cap A = \mathfrak{p}$ and $A/\mathfrak{p}A = \kappa(\mathfrak{p})$. In case
$B = A[X]$, the local ring $B_M$ is a localisation of $B \otimes_A A_{\mathfrak{p}} =
A_{\mathfrak{p}}[X]$, and the fibre ring of $A_{\mathfrak{p}} \to B_M$ is a localisation of
$\kappa(\mathfrak{p})[X]$, hence regular. In case $B = A[[X]]$ then $X \in M$, and
$\mathfrak{p}$ a maximal ideal of $A$, so that $\kappa(\mathfrak{p}) = A/\mathfrak{p}$ and
\[
  B \otimes_A \kappa(\mathfrak{p}) = (A/\mathfrak{p})[[X]] = \kappa(\mathfrak{p})[[X]].
\]
This is a regular local ring, and is the fibre ring of $A_{\mathfrak{p}} \to B_M$. Thus in
either case $B_M$ is Gorenstein by the previous theorem.
\end{proof}

\begin{theorem}\label{thm:23.6}
Let $A$ be a Gorenstein ring containing a field $k$; then for any finitely
generated field extension $K$ of $k$, the ring $A \otimes_k K$ is Gorenstein.
\end{theorem}

\begin{proof}
We need only consider the case that $K$ is generated over $k$ by one element $x$.
If $x$ is transcendental over $k$ then $A \otimes_k K$ is isomorphic to a
localisation of $A \otimes_k k[X] = A[X]$, and since $A[X]$ is Gorenstein, so is
$A \otimes_k K$. If $x$ is algebraic over $k$ then since $K \cong k[X]/(f(X))$ with
$f(X) \in k[X]$ a manic polynomial, we have
\[
  A \otimes_k K = A[X]/(f(X));
\]
now $A[X]$ is Gorenstein and $f(X)$ is a non-zero-divisor of $A[X]$, so that we see
that $A \otimes_k K$ is also Gorenstein.
\end{proof}

\begin{remark}
Theorems 5 and 6 also hold on replacing Gorenstein by Cohen-Macaulay; the
proofs are exactly the same. For complete intersection rings the counterpart of
Theorem 4 also holds, so that the analogs of Theorems 5 and 6 follow; the proof
involves Andr\'e homology (Avramov [1]). As we see in the next theorem, a
slightly weaker form of the same result holds for regular rings.
\end{remark}

\begin{theorem}\label{thm:23.7}
Let $(A, m, k)$ and $(B, n, k')$ be Noetherian local rings, and $A \to B$ a local
homomorphism; set $F = B/mB$. We assume that $B$ is flat over $A$.
\begin{enumerate}[(i)]
  \item If $B$ is regular then so is $A$.
  \item If $A$ and $F$ are regular then so is $B$.
\end{enumerate}
\end{theorem}

\begin{proof}
(i) We have $\Tor^A_i(k, k) \otimes_A B = \Tor^B_i(B \otimes_A k, B \otimes_A k)$, and the right-hand
side is zero for $i > \dim B$. Since $B$ is faithfully flat over $A$, we have
$\Tor^A_i(k, k) = 0$ for $i \gg 0$, so that by \S19, Lemma~\ref{lem:19.1}, (i),
$\projdim_A k < \infty$ and since $\projdim k = \gldim A$, by
Theorem~\ref{thm:19.2}, $A$ is regular.

(ii) Set $r = \dim A$ and $s = \dim F$. Let $\{x_1, \ldots, x_r\}$ be a regular system of
parameters of $A$, and $\{y_1, \ldots, y_s\}$ a subset of $n$ which maps to a regular
system of parameters of $F$. Since $A \to B$ is injective, we can view $A$ as a
subring $A \subset B$. Then $\{x_1, \ldots, x_r, y_1, \ldots, y_s\}$ generates $n$ but
$\dim B = r + s$, so that $B$ is regular.
\end{proof}

\begin{remark}
In Theorem 7, even if $B$ is regular, $F$ need not be. For example, let $k$ be a
field, $x$ an indeterminate over $k$, and $B = k[x]_{(x)}$, $A = k[x^2]_{(x^2)}
\subset B$; then $F = B/x^2B = k[x]/(x^2)$ has a nilpotent element. By Theorem 1,
or directly, we see that $B$ is flat over $A$. (From a geometrical point of view,
this example corresponds to the projection of the plane curve $y = x^2$ onto the
$y$-axis, and, not surprisingly, the fibre over the origin is singular.)
\end{remark}

Consider the following conditions $(R_i)$ and $(S_i)$ for $i = 0, 1, 2, \ldots$ on a
Noetherian ring $A$:
\[
  (R_i) \quad A_P \text{ is regular for all } P \in \Spec A \text{ with } \htop P \le i;
\]
\[
  (S_i) \quad \depth A_P \ge \min(\htop P, i) \text{ for all } P \in \Spec A.
\]
$(S_0)$ always holds. $(S_1)$ says that all the associated primes of $A$ are minimal,
that is $A$ does not have embedded associated primes. $(R_0) + (S_1)$ is the
necessary and sufficient condition for $A$ to be reduced. $(S_i)$ for all $i \ge 0$
is just the definition of a CM ring.

For an integral domain $A$, $(S_1)$ is equivalent to the condition that every prime
divisor of a non-zero principal ideal has height $1$. The characterisation of
normal integral domain given in the corollary to Theorem~\ref{thm:11.5} can be
somewhat generalised as follows.

\begin{theorem}[Serre]\label{thm:23.8}
$(R_1) + (S_2)$ are necessary and sufficient conditions for a Noetherian ring $A$ to
be normal.
\end{theorem}

\begin{proof}
We defined a normal ring (see \S9), by the condition that the localisation at
every prime is an integrally closed domain. The conditions $(R_i)$ and $(S_i)$ are
also conditions on localisations, so that we can assume that $A$ is local.

\textit{Necessity}. This follows from Theorems~\ref{thm:11.2} and \ref{thm:11.5}.

\textit{Sufficiency}. Since $A$ satisfies $(R_0)$ and $(S_1)$ it is reduced, and the
shortest primary decomposition of $(0)$ is $(0) = P_1 \cap \cdots \cap P_r$, where $P_i$
are the minimal primes of $A$. Thus if we set $K$ for the total ring of fractions
of $A$, we have
\[
  K = K_1 \times \cdots \times K_r, \quad \text{with } K_i \text{ the field of fractions of } A/P_i.
\]
First of all we show that $A$ is integrally closed in $K$. Suppose that we have a
relation in $K$ of the form
\[
  (a/b)^n + c_1(a/b)^{n-1} + \cdots + c_n = 0,
\]
with $a, b, c_1, \ldots, c_n \in A$ and $b$ an $A$-regular element. This is equivalent to
a relation
\[
  a^n + \sum_{i=1}^n c_i a^{n-i}b^i = 0
\]
in $A$. Let $P \in \Spec A$ be such that $\htop P = 1$; then by $(R_1)$, $A_P$ is regular,
and therefore normal, so that $a_P \in b_P A_P$, where we write $a_P$, $b_P$ for the
images in $A_P$ of $a$, $b$. Now $b$ is $A$-regular, so that by $(S_1)$, all the prime
divisors of the principal ideal $bA$ have height $1$; thus if
$bA = \mathfrak{q}_1 \cap \cdots \cap \mathfrak{q}_t$ is a shortest primary decomposition and we set
$\mathfrak{p}_i$ for the prime divisor of $\mathfrak{q}_i$, then $a \in bA_{\mathfrak{p}_i} \cap A = \mathfrak{q}_i$ for all
$i$, and hence $a \in bA$, so that $a/b \in A$. Therefore $A$ is integrally closed in
$K$; in particular, the idempotents $e_i$ of $K$, which satisfy $e_i^2 - e_i = 0$, must
belong to $A$, so that from $1 = \sum e_i$ and $e_ie_j = 0$ for $i \neq j$ we get
\[
  A = Ae_1 \times \cdots \times Ae_r.
\]
Now since $A$ is supposed to be local, we must have $r = 1$, so that $A$ is an
integrally closed domain.
\end{proof}

\begin{theorem}\label{thm:23.9}
Let $(A, m)$ and $(B, n)$ be Noetherian local rings and $A \to B$ a local
homomorphism. Suppose that $B$ is flat over $A$, and that $i \ge 0$ is a given
integer. Then
\begin{enumerate}[(i)]
  \item if $B$ satisfies $(R_i)$, so does $A$;
  \item if both $A$ and the fibre ring $B \otimes_A \kappa(\mathfrak{p})$ over every prime ideal
        $\mathfrak{p}$ of $A$ satisfy $(R_i)$, so does $B$.
  \item The above two statements also hold with $(S_i)$ in place of $(R_i)$.
\end{enumerate}
\end{theorem}

\begin{proof}
(i) For $\mathfrak{p} \in \Spec A$, since $B$ is faithfully flat over $A$, there is a prime
ideal of $B$ lying over $\mathfrak{p}$; if we let $P$ be a minimal element among these then
$\htop(P/\mathfrak{p}B) = 0$, so that $\htop P = \htop \mathfrak{p}$. Hence
$\htop \mathfrak{p} \le i \Rightarrow B_P$ is regular, so that by
Theorem~\hyperref[thm:23.7]{7}, $A_{\mathfrak{p}}$ is regular. Also, by the corollary to
Theorem~\hyperref[thm:23.3]{3}, $\depth B_P = \depth A_{\mathfrak{p}}$, so that one sees easily that $(S_i)$ for
$B$ implies $(S_i)$ for $A$.

(ii) Let $P \in \Spec B$ and set $P \cap A = \mathfrak{p}$. If $\htop P \le i$ then we have
$\htop \mathfrak{p} \le i$ and $\htop(P/\mathfrak{p}B) \le i$, hence $A_{\mathfrak{p}}$ and
$B \otimes_A \kappa(\mathfrak{p})$ are both regular, so by Theorem~\hyperref[thm:23.7]{7}, (ii), $B_P$ is regular.
Hence $B$ satisfies $(R_i)$. Moreover, for $(S_i)$ we have
\begin{align*}
  \depth B_P &= \depth A_{\mathfrak{p}} + \depth B \otimes_A \kappa(\mathfrak{p}) \\
             &\ge \min(\htop \mathfrak{p}, i) + \min(\htop P/\mathfrak{p}B, i) \\
             &\ge \min(\htop \mathfrak{p} + \htop P/\mathfrak{p}B, i) = \min(\htop P, i).
\end{align*}
\end{proof}

\begin{corollary*}
Under the same assumptions as Theorem~\hyperref[thm:23.9]{9}, we have
\begin{enumerate}[(i)]
  \item if $B$ is normal (or reduced) then so is $A$;
  \item if both $A$ and the fibre rings of $A \to B$ are normal (or reduced) then
        so is $B$.
\end{enumerate}
\end{corollary*}

\begin{remark}
If $A$ and the closed fibre ring $F = B/mB$ only are normal, then $B$ does not
have to be; for instance, there are known examples of normal Noetherian rings
for which the completion is not normal.
\end{remark}

Finally, we would like to draw the reader's attention to the following obvious,
but useful, fact concerning the fibre ring. Let $\varphi'\colon A' \to B'$ be a ring
homomorphism and $I$ an ideal of $A'$; we set $A = A'/I$, $B = B'/IB'$, and write
$\varphi\colon A \to B$ for the map induced by $\varphi'$. If $\mathfrak{p}' \in \Spec A'$ is such that
$I \subset \mathfrak{p}'$, we set $\mathfrak{p} = \mathfrak{p}'/I$; then the fibre of $\varphi'$ over $\mathfrak{p}'$ coincides
with the fibre of $\varphi$ over $\mathfrak{p}$. To see this,
\[
  B' \otimes_{A'} \kappa(\mathfrak{p}') = B' \otimes_{A'} (A'/\mathfrak{p}')_{\mathfrak{p}'} = B \otimes_A (A/\mathfrak{p})_{\mathfrak{p}} = B \otimes_A \kappa(\mathfrak{p}).
\]
It follows from this that if all the fibre rings of $\varphi'$ have a good property,
the same is true of $\varphi$. For an example of this, see Ex.~\ref{ex:23.2-ex}.

% ============ 习题 ============
\section*{Exercises to \S\ref{sec:23}}
\emph{Prove the following propositions.}

\begin{exercise}\label{ex:23.1-ex}
If $A$ is a Gorenstein local ring then all the fibre rings of $A \to \hat{A}$ are again
Gorenstein; the same thing holds for Cohen-Macaulay.
\end{exercise}

\begin{exercise}\label{ex:23.2-ex}
If $A$ is a quotient of a CM local ring, and satisfies $(S_i)$, then the completion
$\hat{A}$ also satisfies $(S_i)$. In particular, if $A$ does not have embedded
associated primes then neither does $\hat{A}$.
\end{exercise}

\begin{exercise}\label{ex:23.3-ex}
Give another proof of Theorem~\hyperref[thm:23.4]{4} along the following lines:
\begin{enumerate}[(1)]
  \item Using $\Ext^i_A(A/m, A) \otimes_A \hat{B} = \Ext^i_B(F, B)$, show that $B$ Gorenstein
        implies $A$ Gorenstein.
  \item Assuming that $A$ is Gorenstein, prove that $F$ is Gorenstein if and only
        if $B$ is. Firstly reduce to the case $\dim A = 0$. Then prove that
        $\Ext^i_B(F, B) = 0$ for $i > 0$ and $\cong F$ for $i = 0$, and deduce that if
        $0 \to B \to I^{\bullet}$ is an injective resolution of $B$ as a $B$-module then
        $0 \to F \to \Hom_B(F, I^{\bullet})$ is an injective resolution of $F$ as an $F$-module,
        so that, writing $k$ for the residue field of $B$, we have
        $\Ext^i_B(k, B) = \Ext^i_F(k, F)$ for all $i$.
\end{enumerate}
\end{exercise}
